\section{Actual varying-line numerical FS functions are continuous and curve-weak psh}
\label{sec:native-numerical-fs-continuity}

\begin{theorem}[Every actual original rational-line datum has continuous entire evaluation]
\label{thm:native-numerical-fs-data-continuous}
\lean{EnvelopeProperty.NativeNumericalFubiniStudy.FSData.continuous_eval}
\leanok
\uses{def:native-numerical-whole-section-fs-data,thm:scheme-original-global-section-norm-continuous}
On an actual integral scheme projective over its original field, every
actual whole-section numerical FS datum has continuous evaluation on the
same entire original native point space. The original line can be
arbitrary, all weights are real, common zeroes are permitted and the
family may even be zero. No model, ample representative, norm continuity
or section comparison certificate is an input.
\end{theorem}
\begin{proof}
The true norm of each actual original whole section is continuous by the
genuine arbitrary-line norm theorem. The inclusion into extended
nonnegative reals and the extended logarithm are continuous, including at
zero norm. Adding the actual finite real weight preserves continuity.
The finite supremum of these actual continuous extended-real terms is
continuous. The denominator is a positive integer, so multiplication by
its finite nonzero reciprocal is continuous, also at negative infinity.
Their composition is precisely the actual datum's entire original
evaluation.
\end{proof}

\begin{theorem}[Every genuine numerical FS function is continuous]
\label{thm:native-numerical-fs-continuous}
\lean{EnvelopeProperty.NativeNumericalFubiniStudy.IsFS.continuous,EnvelopeProperty.NativeNumericalFubiniStudy.IsRationalFS.continuous}
\leanok
\uses{def:native-numerical-whole-section-fs,thm:native-numerical-fs-data-continuous}
Every numerical FS function of any actual original numerical class is
continuous on the entire original native space. This applies equally to
the actual rational-weight subcollection.
\end{theorem}
\begin{proof}
Choose the actual original defining whole datum. Its entire evaluation
is continuous by the preceding theorem and is the same original
function by the defining equality. The rational case uses the identical
datum after forgetting its extra weight-rationality property.
\end{proof}

\begin{theorem}[Actual numerical FS is contained in its genuine curve-weak psh class]
\label{thm:native-numerical-fs-is-curve-weak-psh}
\lean{EnvelopeProperty.NativeNumericalFubiniStudy.IsFS.isCurveWeakPSH,EnvelopeProperty.NativeNumericalFubiniStudy.IsRationalFS.isCurveWeakPSH,EnvelopeProperty.NativeNumericalFubiniStudy.FSData.isCurveWeakPSH}
\leanok
\uses{thm:native-numerical-fs-continuous,thm:native-numerical-fs-generically-finite,def:native-numerical-curve-weak-psh}
Every genuine numerical FS function of an actual class $\theta$ is
curve-weak numerical psh of that same actual class and function. In
particular, every actual nonzero whole original rational-line datum,
with arbitrary real weights and common zeroes allowed, gives a
curve-weak psh function in its true class $m^{-1}[M]$. The same conclusion
holds for actual rational weights.
\end{theorem}
\begin{proof}
Actual original nonzero sections internally give genuine generic
finiteness, and the preceding theorem gives continuity and hence upper
semicontinuity. Use the constant net consisting of this same original
datum, indexed by the singleton directed preorder. Its actual numerical
classes are constantly $\theta$ and therefore converge to $\theta$ in
the actual curve-weak topology. Its evaluations form a decreasing net,
and the pointwise infimum is the same original evaluation. All defining
conditions are now derived, proving the inclusion without adding an
upper-semicontinuity or nontriviality premise. Apply the same construction
to the rational subcollection, or directly to any given nonzero actual
datum.
\end{proof}

\begin{remark}
This proves the FS inclusion for the actual varying-line, curve-weak
numerical psh definition. It does not assert the general decreasing-psh-net
closure theorem, ample-perturbation regularization, numerical finite
dimensionality, agreement with the paper's Euclidean topology or the
full original/model PSH comparison. The older fixed-ray predicate remains
a distinct scoped endpoint.
\end{remark}
